% ============================================================================
%  chapter10.tex
%  Time, Timers, and I/O Multiplexing
% ============================================================================

\chapter{Time, Timers, and I/O Multiplexing}

\section{The Problem: One Process, Many Channels}

Everything so far treated I/O as a rendezvous: a pipe blocks a reader
until bytes arrive, a signal blocks a \code{waitpid} until a child
finishes. That model has a hidden cost. A process doing blocking I/O
on one channel is \emph{frozen} for everyone else. If a server reads a
socket and no data comes, it cannot notice that a second socket is
ready, or that a timer has fired, or that a child has died.

\begin{verbatim}
   A single-threaded process with two blocking reads:

        process
       +---------------------+
       |  read(fd1, ...)     |  <-- blocks HERE until fd1 has data
       |                     |       ...fd2 could be bursting with
       |  read(fd2, ...)     |       data the whole time, untouched
       |  (never reached)    |
       +---------------------+

   The fix is I/O MULTIPLEXING: watch many descriptors at once, and be
   told which ones are ready. The process waits for ANY of them:

        process
       +---------------------+
       |  poll([fd1, fd2,     |
       |         timer], ...) |  <-- returns when ANY is ready
       |  -> fd2 is ready     |
       |  read(fd2, ...)      |  <-- now safe: it won't block
       +---------------------+
\end{verbatim}

This chapter builds the machinery of the \emph{event loop}: one
process, many channels, no thread per channel. It is the architecture
of every network server (Chapter~17), every window system, and every
game loop. To build it we need three ingredients: a trustworthy clock,
timers that arrive as file descriptors, and a way to ask the kernel
``which of my descriptors are ready right now?''

\begin{osconcept}
Multiplexing is an OS problem because only the kernel knows the true
readiness state of every descriptor at every instant. Your process
cannot cheaply know whether a pipe has data without the kernel telling
it. So the OS offers a \emph{readiness query}: block me, but block me
on \emph{many} channels at once, and wake me the moment any of them
can make progress. That single primitive is what lets one process serve
ten thousand connections.
\end{osconcept}

\section{Time: The Two Clocks}

Every multiplexed program needs time: timeouts, heartbeats, rate
limits. C reads the clock with \code{clock\_gettime}, which fills a
\code{struct timespec}---two fields: seconds and nanoseconds. The
first design decision is \emph{which clock}. There are two you will
use, and confusing them is a classic bug:

\begin{verbatim}
   CLOCK_REALTIME    the wall clock: 2026-08-03 14:00:00
                     - matches `date`, matches the calendar
                     - CAN JUMP: NTP syncs, an admin sets the time,
                       a leap second happens, the value goes
                       backward or forward in one step

   CLOCK_MONOTONIC   time since boot, ever increasing
                     - never jumps, never goes backward
                     - useless for "what day is it?"
                     - perfect for "how long did this take?"

   Rule: measure DURATIONS with MONOTONIC, show TIMES with REALTIME.
   If you time a request with REALTIME and an admin changes the clock
   mid-request, your "elapsed" time can be negative.
\end{verbatim}

\begin{lstlisting}[caption={Timing an interval with the monotonic clock.}]
#include <stdio.h>
#include <time.h>

/* Return the monotonic clock as a double number of seconds. */
double monotonic_seconds(void) {
    struct timespec ts;

    clock_gettime(CLOCK_MONOTONIC, &ts);
    return (double) ts.tv_sec + (double) ts.tv_nsec / 1e9;
}

int main(void) {
    double start;
    double end;
    double elapsed;

    start = monotonic_seconds();

    /* Do some work (a busy pause). */
    {
        volatile unsigned long x = 0;
        unsigned long i;

        for (i = 0; i < 200000000UL; i++) {
            x = x + 1;
        }
    }

    end = monotonic_seconds();
    elapsed = end - start;

    printf("elapsed: %.6f seconds\n", elapsed);
    return 0;
}
\end{lstlisting}

\begin{verbatim}
   struct timespec, drawn:

     struct timespec { long tv_sec;  /* whole seconds     */
                       long tv_nsec; /* nanoseconds 0..999999999 */ };

     tv_sec = 10, tv_nsec = 500000000   ->  10.5 seconds exactly.

   tv_nsec counts nanoseconds, not milliseconds. 1e9 ns = 1 s.
\end{verbatim}

\begin{osconcept}
The kernel owns the clocks. It maintains them, advances the monotonic
clock on every timer tick, and is the only party that can guarantee
``this clock never jumps.'' The realtime/monotonic split is a kernel
design decision: applications measuring intervals get a stable clock,
while the calendar is free to be corrected by an administrator without
breaking every timeout in the system.
\end{osconcept}

\section{Timers}

A \emph{timer} is the kernel's promise to wake you at a future time.
There are two modern styles, and the difference is instructive:
\code{setitimer} delivers its wakeup as a \emph{signal}
(\code{SIGALRM}, Chapter~6), while \code{timerfd} delivers it as a
\emph{file descriptor} that becomes readable---which means a timer can
be watched by \code{poll} alongside real data. The second is the one
this chapter uses, because it makes timers first-class citizens of the
event loop.

\begin{lstlisting}[caption={A periodic timer as a file descriptor.}]
#include <stdio.h>
#include <stdlib.h>
#include <unistd.h>
#include <sys/timerfd.h>
#include <time.h>

/* Create a timer that expires every interval_sec seconds. */
int make_timerfd(int interval_sec) {
    int tfd;
    struct itimerspec its;

    tfd = timerfd_create(CLOCK_MONOTONIC, 0);
    if (tfd < 0) {
        perror("timerfd_create");
        return -1;
    }

    /* it_value: first expiration; it_interval: how often after that. */
    its.it_value.tv_sec = interval_sec;
    its.it_value.tv_nsec = 0;
    its.it_interval.tv_sec = interval_sec;
    its.it_interval.tv_nsec = 0;

    if (timerfd_settime(tfd, 0, &its, NULL) < 0) {
        perror("timerfd_settime");
        close(tfd);
        return -1;
    }
    return tfd;
}
\end{lstlisting}

When a timerfd expires, it becomes \emph{readable}; reading it returns
the number of expirations that have accumulated since the last read.
The word ``accumulated'' matters: if the event loop was busy and four
periods passed, the read returns 4. Timers never silently drop ticks;
they are counted.

\begin{lstlisting}[caption={Drain the timer's accumulated expirations.}]
#include <stdint.h>

/* Read and discard the timer's expiration count. */
void drain_timer(int tfd) {
    uint64_t expirations;
    ssize_t r;

    r = read(tfd, &expirations, sizeof(expirations));
    if (r < 0 && errno != EAGAIN) {
        perror("read timerfd");
    }
}
\end{lstlisting}

\begin{verbatim}
   signal timer vs file-descriptor timer:

     setitimer()        fires SIGALRM; handler runs (ch.6 rules apply:
                        async-signal-safe, flags only)
     timerfd()          becomes readable; poll() notices it; you read
                        a counter -- no signal, no handler, ordinary
                        I/O rules apply

   For an event loop, timerfd wins: the loop already knows how to
   watch descriptors, so a timer is "just another fd."
\end{verbatim}

\begin{osconcept}
A timer is the kernel scheduling a wakeup on your behalf. The kernel
maintains the timer wheel, checks it on every tick, and makes your
descriptor readable (or raises your signal) at the right moment---all
without your process consuming any CPU while it waits. Timers are
kernel state; \code{timerfd\_settime} merely writes the alarm into
that state.
\end{osconcept}

\section{Non-Blocking I/O}

The event loop has one more requirement. When \code{poll} says ``fd2
is readable,'' the subsequent \code{read} should return immediately---
there is data---so blocking is not a problem there. But some situations
defy that: the kernel reports a descriptor ready for reading, you read,
and nothing arrives; or you want to \emph{write} without waiting for
the buffer to clear. For those, a descriptor can be made
\emph{non-blocking}: any operation that would block instead returns
immediately with \code{-1} and \code{errno = EAGAIN}, ``try again
later.''

\begin{verbatim}
   The O_NONBLOCK contract:

     blocking fd:     read() parks the process until data arrives.
     non-blocking fd: read() returns EAGAIN if nothing is there NOW.

   The deal: your code calls poll(), waits for "readable", then reads
   and is guaranteed not to block. Non-blocking mode is the safety net
   that makes that guarantee hold even when the kernel's guess is
   imperfect.

     write(fd, ...) on a full pipe, blocking: parks you.
     write(fd, ...) on a full pipe, non-blocking: returns EAGAIN.
\end{verbatim}

\begin{lstlisting}[caption={Turning a descriptor non-blocking.}]
#include <stdio.h>
#include <stdlib.h>
#include <fcntl.h>
#include <unistd.h>
#include <errno.h>

/* Make fd non-blocking. Returns 0 on success, -1 on failure. */
int set_nonblocking(int fd) {
    int flags;

    flags = fcntl(fd, F_GETFL, 0);
    if (flags < 0) {
        return -1;
    }
    if (fcntl(fd, F_SETFL, flags | O_NONBLOCK) < 0) {
        return -1;
    }
    return 0;
}
\end{lstlisting}

The read-then-check pattern becomes the idiom of every event loop:

\begin{verbatim}
   got = read(fd, buf, sizeof(buf));
   if (got < 0) {
       if (errno == EAGAIN) {
           /* no data right now: go back to poll() and wait */
           continue;
       }
       perror("read");   /* real error */
       break;
   }
   if (got == 0) {
       /* EOF: the writer closed its end */
       break;
   }
   handle(buf, got);
\end{verbatim}

\begin{osconcept}
Blocking is the kernel parking a process; non-blocking is the kernel
\emph{refusing to park it}. The decision is per-descriptor state
(\code{O\_NONBLOCK}) that the kernel honors in every syscall. EAGAIN
is the kernel's answer to ``I won't wait for you''---it gives the
program the choice of what to do with its time instead of deciding
for it. The OS provides both modes because both are essential: blocking
for simple programs, non-blocking for multiplexed ones.
\end{osconcept}

\section{select: The Classic Multiplexer}

\code{select} is the original way to watch many descriptors. You build
a set of descriptors with \code{fd\_set}, tell \code{select} how long
to wait, and it blocks until any of them is ready---then it tells you,
by modifying your set in place, which ones.

\begin{lstlisting}[caption={Watching stdin and a pipe with select.}]
#include <stdio.h>
#include <stdlib.h>
#include <unistd.h>
#include <sys/select.h>

int main(void) {
    int pfd[2];
    fd_set readfds;
    int r;
    int i;

    if (pipe(pfd) < 0) {
        perror("pipe");
        return 1;
    }

    /* fd 0 = stdin, pfd[0] = pipe read end. */
    FD_ZERO(&readfds);
    FD_SET(0, &readfds);
    FD_SET(pfd[0], &readfds);

    /* Watch the highest fd + 1, with no timeout. */
    r = select(pfd[0] + 1, &readfds, NULL, NULL, NULL);
    if (r < 0) {
        perror("select");
        return 1;
    }
    if (r == 0) {
        printf("timeout (unreachable here: we passed no timeout)\n");
        return 1;
    }

    /* select rewrote readfds: only the ready fds remain set. */
    for (i = 0; i <= pfd[0]; i++) {
        if (FD_ISSET(i, &readfds)) {
            printf("fd %d is ready to read\n", i);
        }
    }

    return 0;
}
\end{lstlisting}

\begin{verbatim}
   select()'s shape, and its two famous limits:

     select(max_fd + 1, &readfds, &writefds, &exceptfds, timeout)
            ^ the HIGHEST fd + 1: select scans 0..max_fd every call
            ^ fd_set is a fixed-size bit array, historically 1024
              bits (FD_SETSIZE): fds beyond that are undetectable
            ^ it MODIFIES your sets: you must rebuild them each loop

   select was a breakthrough and remains everywhere, but its fixed
   size and linear scan make it the wrong tool for thousands of fds.
\end{verbatim}

\begin{osconcept}
\code{select} is the kernel performing a \emph{scan}: it checks every
descriptor from 0 to \code{max\_fd}, parks you if none are ready, and
wakes you the moment any is. The \code{fd\_set} is kernel-shared
state that the syscall both reads (what you want) and writes (what is
ready). The design is simple, and its limits---fixed bit width, linear
work per call---are the kernel's data structures showing through.
\end{osconcept}

\section{poll: An Array of Interest}

\code{poll} fixes select's two headaches. Instead of a fixed-size bit
set, you pass an \emph{array of structs}, one per descriptor, each
carrying the events you care about (\code{events}) and receiving the
events that actually happened (\code{revents}). There is no
\code{FD\_SETSIZE} ceiling, no rebuilding of sets, and the array does
not get destroyed by the call.

\begin{lstlisting}[caption={poll: an array of descriptors and interests.}]
#include <stdio.h>
#include <stdlib.h>
#include <unistd.h>
#include <poll.h>

int main(void) {
    int pfd[2];
    struct pollfd fds[2];
    int r;
    int i;

    if (pipe(pfd) < 0) {
        perror("pipe");
        return 1;
    }

    fds[0].fd = 0;              /* stdin */
    fds[0].events = POLLIN;     /* I want to know when it's readable */
    fds[1].fd = pfd[0];         /* the pipe read end */
    fds[1].events = POLLIN;

    /* Block until any watched descriptor is ready, no timeout. */
    r = poll(fds, 2, -1);
    if (r < 0) {
        perror("poll");
        return 1;
    }

    /* r is the number of fds with revents != 0. */
    printf("%d descriptor(s) ready:\n", r);
    for (i = 0; i < 2; i++) {
        if (fds[i].revents & POLLIN) {
            printf("  fd %d is readable\n", fds[i].fd);
        }
        if (fds[i].revents & POLLERR) {
            printf("  fd %d had an error\n", fds[i].fd);
        }
        if (fds[i].revents & POLLHUP) {
            printf("  fd %d hung up (peer closed)\n", fds[i].fd);
        }
    }

    return 0;
}
\end{lstlisting}

\begin{verbatim}
   struct pollfd, drawn:

     struct pollfd {
         int fd;          /* which descriptor to watch */
         short events;    /* what I'm interested in: POLLIN, POLLOUT */
         short revents;   /* what actually happened (kernel writes)  */
     };

   - POLLIN   data can be read without blocking
   - POLLOUT  space to write without blocking
   - POLLHUP  the peer closed (a hangup)
   - POLLERR  an error condition

   events stays untouched; revents is filled by the kernel. The array
   is rebuilt only when your interests change, not every loop.
\end{verbatim}

poll still tells the kernel your full interest list on every call, and
the kernel scans it each time. For a few dozen descriptors that is
irrelevant; for ten thousand it dominates. That is the regime
\code{epoll} was built for.

\section{epoll: The Scalable Multiplexer}

Linux's answer to scaling is \code{epoll}. The core idea is a
\emph{persistent interest table in the kernel}. You register each
descriptor once, with \code{epoll\_ctl}, and the kernel watches it
forever---no need to resend your list on every wait. When descriptors
become ready, \code{epoll\_wait} hands back only the ones that are,
so the work per call is proportional to the number of \emph{ready}
descriptors, not the total watched.

\begin{lstlisting}[caption={epoll: register once, wait many times.}]
#include <stdio.h>
#include <stdlib.h>
#include <unistd.h>
#include <sys/epoll.h>

int main(void) {
    int pfd[2];
    int epfd;
    struct epoll_event ev;
    struct epoll_event ready[8];
    int r;
    int i;

    if (pipe(pfd) < 0) {
        perror("pipe");
        return 1;
    }

    epfd = epoll_create1(0);
    if (epfd < 0) {
        perror("epoll_create1");
        return 1;
    }

    /* Register the pipe read end, and stdin, for readability. */
    ev.events = EPOLLIN;
    ev.data.fd = pfd[0];
    epoll_ctl(epfd, EPOLL_CTL_ADD, pfd[0], &ev);

    ev.events = EPOLLIN;
    ev.data.fd = 0;
    epoll_ctl(epfd, EPOLL_CTL_ADD, 0, &ev);

    /* Wait. epoll_wait returns only the fds that are actually ready. */
    r = epoll_wait(epfd, ready, 8, -1);
    if (r < 0) {
        perror("epoll_wait");
        return 1;
    }

    for (i = 0; i < r; i++) {
        if (ready[i].events & EPOLLIN) {
            printf("fd %d is readable\n", ready[i].data.fd);
        }
    }

    close(epfd);
    return 0;
}
\end{lstlisting}

\begin{verbatim}
   The epoll lifecycle:

     epoll_create1()      make an empty interest table, get its fd
     epoll_ctl(ADD)       "start watching fd X"   (register once)
     epoll_wait()         "block until some are ready, tell me which"
     epoll_ctl(DEL)       "stop watching fd X"

   Why it scales: the kernel keeps your list; epoll_wait() walks only
   the ready queue. With 10,000 idle connections and 10 active ones,
   each wait returns 10 events -- not a 10,000-entry scan.

   select/poll:  O(total fds)  per call   (rescan everything)
   epoll:        O(ready fds)  per call   (hand back only the active)
\end{verbatim}

\begin{osconcept}
epoll moves the \emph{state}---your interest list---from user memory
into the kernel, and that move is the whole optimization. The kernel
can associate each descriptor with the process watching it, wake the
process only when something in \emph{its} table changes, and hand back
a short list instead of re-scanning the world. Persistent, kernel-owned
interest tables are how the OS lets one process scale to tens of
thousands of simultaneous channels.
\end{osconcept}

\section{The Event Loop: A Reactor}

We now have every ingredient. The \emph{event loop} (also called a
\emph{reactor}) is the pattern that ties them together, and it is the
heart of every serious server:

\begin{verbatim}
   The reactor pattern, one diagram:

        +----------------------------------------------------+
        |                THE EVENT LOOP                      |
        |                                                    |
        |   +--------+   wait on ALL fds with a timeout      |
        |   | poll/  |   (including timerfds)                |
        |   | epoll  |                                       |
        |   | wait   |                                       |
        |   +--------+                                       |
        |       |                                            |
        |       v                                            |
        |   for each READY fd:                               |
        |       if timer:     drain it, do periodic work     |
        |       if data:      read (non-blocking), handle    |
        |       if EOF/hup:   close, forget this fd          |
        |       ^                                            |
        |       +--------- loop forever --------------------+
        +----------------------------------------------------+

   Invariant: no blocking call exists inside the loop except the wait
   itself. Every read/write is attempted only when the descriptor is
   known ready, and even then non-blockingly.
\end{verbatim}

The rules that keep a reactor healthy are few and absolute:

\begin{itemize}
  \item \emph{Never block in the handler.} If handling one fd could
    block, the whole world stops. Work that might stall is deferred,
    queued, or done elsewhere.
  \item \emph{Non-blocking everywhere.} A read that returns
    \code{EAGAIN} is not an error; it is the loop telling you to go
    wait again.
  \item \emph{Timers are fds.} Heartbeats and timeouts ride the same
    wait as data, so the loop wakes at deadlines even when no I/O
    arrives.
\end{itemize}

\section{Practical: A Multi-Channel Event Loop}

Let us build the real thing: one process watches \emph{two pipes and a
timer} with \code{poll}, and reacts to whichever is ready---exactly
the shape a network server's main loop will take in Chapter~17. Two
children write at different times; the loop notices each one when it
arrives, and ticks on the timer meanwhile.

\begin{lstlisting}[caption={mux.c: one loop, many channels, no thread.}]
#include <stdio.h>
#include <stdlib.h>
#include <unistd.h>
#include <poll.h>
#include <sys/timerfd.h>
#include <time.h>
#include <errno.h>
#include <stdint.h>

int make_timerfd(int interval_sec) {
    int tfd;
    struct itimerspec its;

    tfd = timerfd_create(CLOCK_MONOTONIC, 0);
    if (tfd < 0) {
        return -1;
    }
    its.it_value.tv_sec = interval_sec;
    its.it_value.tv_nsec = 0;
    its.it_interval.tv_sec = interval_sec;
    its.it_interval.tv_nsec = 0;
    if (timerfd_settime(tfd, 0, &its, NULL) < 0) {
        close(tfd);
        return -1;
    }
    return tfd;
}

void drain_timer(int tfd) {
    uint64_t expirations;
    ssize_t r;

    r = read(tfd, &expirations, sizeof(expirations));
    if (r < 0 && errno != EAGAIN) {
        perror("read timerfd");
    }
}

int main(void) {
    int pfd_a[2];
    int pfd_b[2];
    int tfd;
    struct pollfd fds[3];
    int remaining;
    int r;
    int i;

    if (pipe(pfd_a) < 0 || pipe(pfd_b) < 0) {
        perror("pipe");
        return 1;
    }

    /* Two writer children, each producing one line after a delay. */
    if (fork() == 0) {
        sleep(2);
        write(pfd_a[1], "from A\n", 7);
        _exit(0);
    }
    if (fork() == 0) {
        sleep(4);
        write(pfd_b[1], "from B\n", 7);
        _exit(0);
    }

    /* The parent never writes; drop both write ends. */
    close(pfd_a[1]);
    close(pfd_b[1]);

    tfd = make_timerfd(1);
    if (tfd < 0) {
        return 1;
    }

    /* Register all three channels. */
    fds[0].fd = pfd_a[0];
    fds[0].events = POLLIN;
    fds[1].fd = pfd_b[0];
    fds[1].events = POLLIN;
    fds[2].fd = tfd;
    fds[2].events = POLLIN;

    remaining = 2;          /* the two pipes we must hear from */

    while (remaining > 0) {
        r = poll(fds, 3, -1);
        if (r < 0) {
            if (errno == EINTR) {
                continue;              /* retry after a signal */
            }
            perror("poll");
            return 1;
        }

        for (i = 0; i < 3; i++) {
            char buf[64];
            ssize_t got;

            if ((fds[i].revents & POLLIN) == 0) {
                continue;
            }

            if (fds[i].fd == tfd) {
                drain_timer(tfd);
                printf("tick\n");
                continue;
            }

            /* A pipe is readable: read without blocking. */
            got = read(fds[i].fd, buf, sizeof(buf));
            if (got > 0) {
                printf("fd %d says: %.*s", fds[i].fd, (int) got, buf);
            } else {
                /* EOF: the writer closed. Stop watching this pipe. */
                printf("fd %d closed\n", fds[i].fd);
                fds[i].fd = -1;        /* poll ignores negative fds */
                remaining = remaining - 1;
            }
        }
    }

    printf("both channels done; exiting\n");
    return 0;
}
\end{lstlisting}

\begin{verbatim}
   Running it:

     $ ./mux
     tick                 (t=1: timer fires first)
     tick                 (t=2: timer again)
     fd 4 says: from A    (t=2: child A's line arrives the same second)
     tick                 (t=3)
     tick                 (t=4: timer)
     fd 5 says: from B    (t=4: child B's line)
     fd 5 closed          (writer B exited)
     fd 4 closed          (writer A's EOF was noticed too)
     both channels done; exiting

   One process, three channels, zero threads, zero busy-waiting. The
   kernel did all the waiting; the loop just dispatched what was ready.
\end{verbatim}

\begin{researchnote}[Event-driven servers at scale]
The event loop you just built is not a textbook toy; it is the exact
architecture that made early web servers fast. In 1999, researchers at
Rice University published ``Flash,'' a web server built around
event-driven I/O---a single loop multiplexing thousands of
connections---and showed it outperforming the threaded servers of the
day by wide margins.

\begin{verbatim}
   What Flash demonstrated (from the paper's own claims):

       "An event-driven server can handle tens of thousands of
        concurrent connections on commodity hardware, and its
        scalability comes from avoiding one-thread-per-connection
        and from careful handling of each I/O operation."

   In this chapter's terms: Flash used exactly the pattern of Section
   9.7 -- one poll-style wait, non-blocking handlers, no thread per
   connection. The lesson has aged well; today's event loops
   (epoll, kqueue, io_uring) are its descendants.
\end{verbatim}

When you build the scalable server of Chapter~17, you will be
repeating this paper's architecture on modern primitives---and now you
know where the idea came from.

\paperref{V.~S.~Pai, P.~Druschel, W.~Zwaenepoel, ``Flash: An Efficient
and Portable Web Server,'' USENIX Annual Technical Conference,
1999}{https://www.usenix.org/legacy/events/usenix99/full_papers/pai/pai.pdf}
\end{researchnote}

\section{Structures in Bytes}

Everything the event loop touches has a small, fixed byte shape:

\begin{verbatim}
   struct timespec        16 bytes   offset size field
                                     0      8    tv_sec
                                     8      8    tv_nsec
   struct itimerspec      32 bytes   0..15  it_interval
                                     16..31 it_value
   struct pollfd           8 bytes   0      4    fd
                                     4      2    events
                                     6      2    revents
   struct epoll_event     12 bytes   0      4    events (EPOLLIN, ...)
                                     4      8    data   (a 64-bit cookie)
   fd_set (select)       128 bytes   one bit per descriptor (1024 bits)

   a timerfd's readable payload: 8 bytes = a uint64_t expiration count
\end{verbatim}

Read the shapes and the API is predictable: \code{poll} takes an array
of 8-byte \code{pollfd} records and the kernel writes 2 bytes of
\code{revents} into each; \code{epoll\_wait} returns an array of
12-byte \code{epoll\_event} records; \code{clock\_gettime} fills 16
bytes. These layouts are fixed by the ABI---which is exactly why you
should never read or write them as raw bytes yourself; the helper
calls exist to keep you from doing so.

\section{Summary}

\begin{verbatim}
   The multiplexing toolbox, complete:

   clock_gettime(CLOCK_MONOTONIC)  interval timing that never jumps
   timerfd_create()                a timer as a readable descriptor
   fcntl(F_SETFL, O_NONBLOCK)      read without ever parking
   EAGAIN                          "nothing right now, wait again"
   select()                        the classic bit-set multiplexer
   poll()                          array of pollfd (events/revents)
   epoll()                         kernel-held interest table
   the event loop                  one wait, dispatch what's ready

   Rules that keep an event loop correct:
     1. no blocking calls inside handlers;
     2. reads may return EAGAIN -- go back to the wait;
     3. EOF (got == 0) and POLLHUP mean close and forget;
     4. timers are fds: wait on them like data;
     5. use MONOTONIC for durations, REALTIME for display.
\end{verbatim}

\begin{osconcept}
Multiplexing is the OS delivering on a hard promise: one process, many
channels, and the kernel does the waiting. It is an OS problem because
readiness is kernel knowledge and parking is kernel power. The
evolution from select to epoll is the OS learning to hold your
interest lists, so that scale no longer means more threads---just a
better wait.
\end{osconcept}

We have seen one process do the work of many. The next chapter goes
the other way: instead of multiplexing one process across channels, we
give a process \emph{many threads of execution} inside one address
space---sharing memory, racing on every variable, and forcing us to
learn the rules of concurrency.

